% =====================================================================
\section{Introduction}\label{sec:intro}
% =====================================================================

\subsection{The problem in one paragraph}
In a normal Internet of Vehicles (IoV), roadside units (RSUs) mounted on poles verify every vehicle and relay data to a cloud
Trusted Authority (TA). After an earthquake or flood the poles, the fibre and the power are the first things to fail. Rescue convoys
and evacuation traffic then need secure coordination most, and nobody is left to verify anyone. Tan, Zheng and Vijayakumar~\cite{base}
call this setting the \emph{infrastructure-less IoV}. Their answer is a \emph{tethered UAV} (TUAV), a drone powered by a cable from an
emergency communications response vehicle (ECRV), which hovers for hours and acts as a flying edge server for a swarm of ordinary UAVs.

\subsection{What the base paper contributes}
The paper~\cite{base} makes three contributions, which we presented in Phase~1:
\begin{enumerate}[label=\textbf{(\roman*)}]
  \item a TUAV-enabled edge architecture that replaces the dead RSU and gives three-dimensional coverage;
  \item a certificateless \emph{group} authentication of $n$ UAVs at once, with the costly bilinear pairings moved to the TA;
  \item a group key distributed with the Chinese Remainder Theorem (CRT) and updated cheaply when UAVs join or leave.
\end{enumerate}
It reports 151\,ms to authenticate 120 UAVs, 104 bytes per request, and all eleven security properties of its Table~II.

\subsection{What we found}
We re-derived every equation of Section~IV and re-did the cost and byte accounting of Section~VI. The scheme is elegant, but it
does not deliver its title property, and several entries in its security table do not hold:
\begin{itemize}
  \item \textbf{It still depends on infrastructure.} The TUAV cannot verify anyone by itself. Steps GA1--GA2 send every request over a
        satellite link to the TA, which computes every pairing. If the link is lost, authentication stops completely (Finding~\ref{f:online}).
  \item \textbf{The TA can be any UAV.} The credential $\Gamma_i$ depends only on $s_i$, which the TA issued and stores, and on public
        values. So the TA can forge any UAV (key escrow), and no UAV message is non-repudiable (Finding~\ref{f:escrow}).
  \item \textbf{Revocation does not revoke.} The published CRT polynomial reveals every member's decrypting key to every other member, so a
        removed UAV still recovers the new group key (Finding~\ref{f:crt}). We confirmed this with a script.
  \item \textbf{The numbers leave out the TA's work.} With the TA's pairings counted, group authentication of 120 UAVs costs about 3720\,ms
        rather than 151\,ms. The key acknowledgements carry about 143\,KiB each, not 104 bytes (Finding~\ref{f:cost}).
  \item \textbf{The proof of Theorem~1 is empty.} Its CDH ``instance'' has a public exponent, so the reduction proves nothing (Finding~\ref{f:proof}).
\end{itemize}

\subsection{Our contributions}
\begin{keybox}
Keep the TUAV architecture, and make the trust root \emph{offline}. The TA signs credentials before deployment and can trace a
misbehaving entity afterwards, but no step in the field ever waits for it. Every UAV, vehicle and TUAV can verify every other one
using only the TA's public key.
\end{keybox}

\begin{enumerate}[label=\textbf{C\arabic*.},leftmargin=*]
  \item[\textbf{A.}] \textbf{Analysis of the base scheme}: five findings, each with what fails, where, why and its impact (Section~\ref{sec:findings}).
  \item[\textbf{C1.}] \textcolor{uavc}{\textbf{Autonomous UAV tier.}} Escrow-free, pairing-free certificateless credentials verified \emph{locally} by the TUAV in a
        small-exponent batch; a bounded fallback when a batch fails; group keys delivered per member under fresh ECDH keys, so revocation is secure and the
        traffic is linear in $n$ (Section~\ref{sec:c1}).
  \item[\textbf{C2.}] \textcolor{vehc}{\textbf{Vehicle tier.}} Vehicles, which the base paper explicitly leaves out, are authenticated under the \emph{same} offline trust root. Their requests
        are relayed by authenticated UAVs, which also filter floods, and they receive a regional key for V2V traffic (Section~\ref{sec:c2}).
  \item[\textbf{C3.}] \textcolor{hoc}{\textbf{TUAV-to-TUAV handover.}} Neighbouring TUAVs share a link key. A member moving between cells presents an encrypted ticket and is
        admitted using symmetric operations plus one ECDH, without a full re-authentication and without the TA (Section~\ref{sec:c3}).
  \item[\textbf{E.}] \textbf{Security and cost analysis}: theorems with proof sketches, an updated security table, and a cost and byte comparison that
        uses the base paper's own constants (Sections~\ref{sec:security}--\ref{sec:perf}).
\end{enumerate}

\subsection{Organisation}
Section~\ref{sec:base} recalls the base scheme. Section~\ref{sec:findings} presents our findings. Section~\ref{sec:model} gives the system model,
threat model and design goals. Section~\ref{sec:prelim} introduces the building blocks and notation. Sections~\ref{sec:setup}--\ref{sec:c3} present the
framework. Section~\ref{sec:security} analyses security, Section~\ref{sec:perf} performance, Section~\ref{sec:discussion} limitations, and
Section~\ref{sec:plan} the implementation and evaluation plan. Section~\ref{sec:conclusion} concludes.

% =====================================================================
\section{The Base Scheme in Brief}\label{sec:base}
% =====================================================================
This section uses the base paper's notation. Figure~\ref{fig:base} shows the message flow we presented in Phase~1.

\begin{figure}[H]
\centering
\resizebox{0.97\linewidth}{!}{%
\begin{tikzpicture}[node distance=1cm]
  \node[ent=deep]  (ta)   at (0,0)   {TA (cloud)};
  \node[ent=hoc]   (tu)   at (6,0)   {TUAV};
  \node[ent=uavc]  (uav)  at (12,0)  {$n$ UAVs};
  \foreach \n in {ta,tu,uav}{ \draw[life] (\n.south) -- ++(0,-8.2); }
  \draw[msg=deep,dashed] (0,-1.0) -- node[lbl,above]{IS1 $\langle id_{tu},s_{tu}\rangle$ (offline)} (6,-1.0);
  \draw[msg=deep,dashed] (0,-1.7) -- node[lbl,above,pos=0.75]{IS2 $\langle id_i,s_i\rangle$ to each UAV (offline)} (12,-1.7);
  \draw[msg=hoc] (6,-2.4) -- node[lbl,above]{IS5 $\langle ts_{tu},ID_{tu},R_{tu},Q_{tu}\rangle$ broadcast} (12,-2.4);
  \draw[msg=uavc] (12,-3.2) -- node[lbl,above]{US3 $\langle Request,ts_2,ID_i,\xi_i,\upsilon_i,\Gamma_i\rangle\times n$} (6,-3.2);
  \draw[msg=bad,very thick] (6,-4.0) -- node[lbl,above]{GA1 upload (satellite)} (0,-4.0);
  \node[act=deep,anchor=east] at (-0.2,-4.75) {GA2: map $ID_i\!\to\! id_i$,\\ $3n$ pairings};
  \draw[msg=bad,very thick] (0,-5.5) -- node[lbl,above]{GA2 $\langle\sigma_i,\psi_i,\mu_i,\delta_i\rangle$ (satellite)} (6,-5.5);
  \node[act=hoc] at (6,-6.25) {GA3: one product equation};
  \draw[msg=hoc] (6,-7.0) -- node[lbl,above]{GD4 $\langle ACK,ts_3,ID_i,S_i,\{a_k\}_{k=0}^{n}\rangle$ to each} (12,-7.0);
  \node[act=uavc,anchor=west] at (12.2,-7.7) {GD5--6: check $S_i$,\\ $\kappa_{tu}=\Lambda(\sigma_i)\bmod\sigma_i$};
  \draw[bad,thick,decorate,decoration={brace,amplitude=5pt}] (-0.1,-3.75) -- (-0.1,-5.75);
\end{tikzpicture}}
\caption{Base scheme~\cite{base}. Red arrows cross the satellite link. GA3 cannot run until GA2 returns.}
\label{fig:base}
\end{figure}

\noindent The equations we use later (verbatim from~\cite{base}, Section~IV):
\begin{align}
  &\text{IS4:} && ID_{tu}=h_2(id_{tu},r_{tu}),\quad R_{tu}=r_{tu}\,h_2(ts_{tu},s_{tu})\in\mathbb{Z}_q,\quad Q_{tu}=s_{tu}\,h_2(ID_{tu},r_{tu})\,P \label{eq:is4}\\
  &\text{US1:} && \xi_i=h_1(r_i),\qquad ID_i=h_3(id_i,ts_2,\xi_i)\\
  &\text{US2:} && \upsilon_i=h_4(ID_i,ID_{tu},ts_2,\xi_i),\qquad
     \Gamma_i=\big[s_i\,h_3(ID_i,ts_2,s_i)\,R_{tu}-h_1(s_i\xi_i)\big]P-\upsilon_i\,h_1(r_i)\,Q_{tu} \label{eq:gamma}\\
  &\text{GA2:} && \sigma_i=h_1(s_i\xi_i)P,\ \ \psi_i=\hat e(\Gamma_i+\sigma_i,P),\ \ \mu_i=\hat e(h_3(ID_i,ts_2,s_i)P,s_iP),\ \ \delta_i=\hat e(s_{tu}P,\upsilon_i\xi_iP)\\
  &\text{GA3:} && \textstyle\prod_{i}\psi_i\,\big(\prod_i\delta_i\big)^{h_2(ID_{tu},r_{tu})}\overset{?}{=}\prod_i\mu_i^{R_{tu}}\label{eq:ga3}\\
  &\text{GD1--2:} && \omega_i=\textstyle\prod_{j\ne i}\sigma_j,\ \ \varphi_i\equiv\omega_i^{-1}\bmod\sigma_i,\ \
     \Lambda(x)=\kappa_{tu}\sum_{i=1}^{n}\omega_i\varphi_i+\prod_{i=1}^{n}(x-\sigma_i)=\sum_{k=0}^{n}a_kx^k \label{eq:lambda}\\
  &\text{GD6:} && \kappa_{tu}=\Lambda(\sigma_i)\bmod\sigma_i
\end{align}
